\clearpage
\VerdierPhysicalPageAnchor{0247}%
\begingroup\small\setlength{\parskip}{0pt}\setlength{\emergencystretch}{1.5em}
\setlength{\abovedisplayskip}{4pt plus 2pt minus 2pt}%
\setlength{\belowdisplayskip}{4pt plus 2pt minus 2pt}%

\VerdierSemanticAnchor{ch4d-4-6-5-resolution-type}
\noindent を備え、入射型分解 (1.1.4) となるものをいう。

\VerdierSemanticAnchor{ch4d-4-6-5-terminology}
したがって、われわれが\emph{複体 $Y$ の Cartan--Eilenberg 入射分解}と
呼ぶものは、これらの著者が単に\emph{複体 $Y$ の入射分解}と呼ぶもの
（[1]、第 XVII 章）であるが、本書では後者の用語を用いない。$Y$ の
Cartan--Eilenberg 入射分解は、複体 $Y$ の各成分と $Y$ の各微分の核の
入射型分解を誘導する［\emph{同所引用文献}］。

\VerdierSemanticAnchor{ch4d-4-6-5-finite-length}
整数 $n_0$ が存在し、$n>n_0$ ならば $I^{\bullet,n}=0$ となるとき、
Cartan--Eilenberg 分解
$\varepsilon\colon Y\longrightarrow I^{\bullet\bullet}$ は
\emph{有限長}であるという。複体 $Y$ が有限長の Cartan--Eilenberg 入射分解を
もつための必要十分条件は、$Y$ のコホモロジー対象と $Y$ の各微分の像の
入射次元が有界であることである。\emph{同所引用文献}で得られた結果を
想起しよう：

\VerdierSemanticAnchor{ch4d-4-6-6-a}
\medskip\noindent\textbf{命題 4.6.6.}\quad
\textit{a) 圏 $\mathcal A$ が十分な入射対象をもつとき (3.1.5)、
$\mathcal A$ の任意の単純複体 $Y$ は Cartan--Eilenberg 入射分解をもつ。}

\VerdierSemanticAnchor{ch4d-4-6-6-b}
\noindent\textit{b)
$f\colon X\longrightarrow Y$ を単純複体の射とし、
$\varepsilon\colon X\longrightarrow I^{\bullet\bullet}$ および
$\varepsilon'\colon Y\longrightarrow J^{\bullet\bullet}$ を二つの
Cartan--Eilenberg 入射分解とする。このとき二重複体の射}

\VerdierSemanticAnchor{ch4d-4-6-6-b-morphism}
\[
\textit{$g\colon I^{\bullet\bullet}\longrightarrow J^{\bullet\bullet}$}
\]
\textit{で $f$ の上にあるもの、すなわち二重複体の射の図式}

\VerdierSemanticAnchor{ch4d-4-6-6-b-diagram}
\[
\begin{tikzcd}[column sep=5.0em,row sep=3.0em]
X \arrow[r,"f"] \arrow[d,"\varepsilon"'] &
Y \arrow[d,"\varepsilon'"] \\
I^{\bullet\bullet} \arrow[r,"g"'] &
J^{\bullet\bullet}
\end{tikzcd}
\]
\textit{を可換にするものが存在する。}

\VerdierSemanticAnchor{ch4d-4-6-6-c}
\noindent\textit{c) $f$ の上にある二つの射
$g,g'\colon I^{\bullet\bullet}\longrightarrow J^{\bullet\bullet}$ は
ホモトピックである（第 I 章、2.5.1 の二重複体の射のホモトピーの意味で）。}

\VerdierSemanticAnchor{ch4d-4-6-6-proof}
証明は\emph{同所引用文献}を参照されたい。

\VerdierSemanticAnchor{ch4d-4-6-7}
\medskip\noindent\textbf{4.6.7.}\quad
$Y$ を $\operatorname{comp}(\mathcal A)$ の対象とし、
$Y\longrightarrow I^{\bullet\bullet}$ を $Y$ の Cartan--Eilenberg 入射分解と
する。二重複体 $I^{\bullet\bullet}$ は、第一次数または第二次数でフィルター
付けすることに応じて、$D(\mathcal A)$ に値をとる二つのスペクトル対象を
生じる (4.1.7)。第二次数でフィルター付けした $I^{\bullet\bullet}$ が定める、
$D(\mathcal A)$ に値をとるスペクトル対象を、\emph{$Y$ の
Cartan--Eilenberg スペクトル対象}という（考えている Cartan--Eilenberg
分解に関して）。
\endgroup
